% Л13 -- Данные: семейства баз и ADO.NET. Slides.
%
% Companion to the spoken script in ../L13-databases-and-ado.md, which is the
% source of truth: the slides carry diagrams, code, real terminal output and
% short marks, and the delivery lives in the script only -- this deck carries
% no speaker notes. Slide text is in Russian because that is the content;
% everything else in this file is English.
%
% Build:  make          (latexmk -pdf; engine: pdflatex)
%
% Every C# listing was compiled on dotnet SDK 10.0 in a scratch
% `dotnet new console -n HelloWorld` with Npgsql 10.0.3; terminal output on
% the slides is captured from runs of that project against PostgreSQL 18 in
% a container. See ../CLAUDE.md for the verification table.
%
% House rules: ../CLAUDE.md and ../../../CLAUDE.md. Theme: ../../../../theme/.

\documentclass[aspectratio=169,10pt,t]{beamer}

\usetheme{dotnet}
\usetikzlibrary{arrows.meta,positioning,calc,decorations.pathreplacing,fit}
\usepackage{array}   % ragged-right table columns

% A ragged-right paragraph column: justified cells in narrow columns hyphenate
% badly, and a split word is the one thing the back row cannot read.
\newcolumntype{L}[1]{>{\raggedright\arraybackslash}p{#1}}

\dotnetlecture{Л13}

\title[Л13. Данные: семейства баз и ADO.NET]{Данные: семейства баз и ADO.NET}
\author{Дмитрий Курочкин}
\date{}

% Code sizes, one step down from \footnotesize. Smaller than 6.6pt does not
% carry to the back of the room.
\newcommand{\codesm}{\ttfamily\fontsize{7.6}{9.0}\selectfont}
\newcommand{\codexs}{\ttfamily\fontsize{7.0}{8.3}\selectfont}
\newcommand{\codexxs}{\ttfamily\fontsize{6.8}{7.8}\selectfont}
\newcommand{\termsm}{\ttfamily\fontsize{6.6}{7.9}\selectfont}

% A mark at the start of a listing line, for labels set beside the listing.
% Used through the listing's escapechar; the labels are drawn by an overlay
% picture after the listing, so the build needs two passes (latexmk does that).
\newcommand{\mk}[1]{\tikz[remember picture,overlay]\coordinate (#1) at (0,0.55ex);}

\tikzset{
  rawbrace/.style={decorate,decoration={brace,amplitude=4pt},hairline,
                   line width=0.8pt},
  rawlab/.style={anchor=west,font=\footnotesize,inner sep=0pt},
  box/.style={draw=hairline,line width=0.8pt,fill=codebg,inner sep=4pt,
              align=center,minimum height=0.9cm},
  arr/.style={-{Straight Barb[length=1.6mm]},muted,line width=0.7pt},
  arra/.style={-{Straight Barb[length=1.8mm]},accent,line width=0.9pt},
  arrl/.style={-{Straight Barb[length=1.8mm]},alt,line width=0.9pt},
}

\begin{document}
\shorthandoff{"}

% ------------------------------------------------------------------ title --
% 1
\begin{frame}[plain]
  \titlepagednc{Лекция 13}{Данные: семейства баз и ADO.NET}{}{Дмитрий Курочкин}
\end{frame}

% ==================================================== 1. DATABASE FAMILIES ==
\section{1. Какие бывают базы}

% 2
\begin{frame}[plain]
  \sectionopener{1}{Какие бывают базы}{Выбор базы \hi{переживёт} весь
    остальной код.}
\end{frame}

% 3
\begin{frame}[plain]
  \vote{Какие базы данных вы знаете? И чем они отличаются?}
\end{frame}

% 4
\begin{frame}{Пять семейств}
  \slidetop
  \claim{Семейство --- это \hi{форма данных} и вопросы к ним.}

  \begin{center}
    \begin{tikzpicture}
      \foreach \i/\t/\s/\e in {
          1/реляционные/{таблицы\\и связи}/{PostgreSQL\\MySQL, SQLite},
          2/документные/{JSON-\\документы}/{MongoDB\\CouchDB},
          3/ключ-значение/{словарь}/{Redis},
          4/колоночные/{столбцы для\\аналитики}/{ClickHouse\\Cassandra},
          5/графовые/{узлы\\и рёбра}/{Neo4j}} {
        \node[box,minimum width=2.45cm,minimum height=1.5cm,inner sep=2pt,
              font=\scriptsize] (f\i) at ({(\i-1)*2.78},0)
          {{\footnotesize\bfseries \t}\\[0.3em]\s};
        \node[font=\scriptsize,text=muted,align=center,below=0.35em] at (f\i.south) {\e};
      }
      \node[box,draw=accent,line width=1.0pt,minimum width=2.45cm,
            minimum height=1.5cm,inner sep=2pt,font=\scriptsize] at (0,0)
        {{\footnotesize\bfseries реляционные}\\[0.3em]таблицы\\и связи};
    \end{tikzpicture}
  \end{center}

  \slidegap
  \aside{«Чем отличаются?» --- «синтаксисом» --- это разница \lo{внутри} одного
    семейства. Не знаете, что нужно, --- нужно \hi{реляционное}.}
\end{frame}

% 5
\begin{frame}{Реляционные и документные}
  \slidetop
  \versus{Реляционные}{%
    \begin{cbul}{accent}
      \item таблицы со строгой схемой
      \item связь --- ссылкой по ключу
      \item SQL и транзакции
      \item любой вопрос поперёк данных
    \end{cbul}}{Документные}{%
    \begin{cbul}{alt}
      \item документ со своей структурой
      \item связь --- вложенностью
      \item читается целиком, одним куском
      \item поперёк документов --- трудно
    \end{cbul}}

  \slidegap
  \hrulethin{0.42}
  \rulegap
  \aside{Граница размыта: PostgreSQL хранит JSON в колонке, с индексом и запросами
    внутрь. Часто это \hi{одна} база вместо двух.}
\end{frame}

% 6
\begin{frame}{Схема есть всегда}
  \slidetop
  \claim{«Без схемы» не значит «без \lo{структуры}».}

  \versus{Схема в базе}{%
    \begin{cbul}{accent}
      \item записана и проверяется
      \item меняется миграцией
      \item старых форм не остаётся
    \end{cbul}}{Схема в коде}{%
    \begin{cbul}{alt}
      \item размазана по тем, кто читает
      \item меняется неявно
      \item код читает обе формы
    \end{cbul}}
\end{frame}

% 7
\begin{frame}{Остальные семейства}
  \slidetop
  {\renewcommand{\arraystretch}{1.5}\small
  \begin{tabular}{@{}L{3.5cm}L{5.6cm}L{3.7cm}@{}}
    & \colhead{accent}{быстро} & \colhead{alt}{не для этого} \\
    \textbf{ключ-значение}      & кэш, сессии, счётчики
      & поиск по данным \\
    \textbf{аналитические}      & сумма по миллиардам строк
      & один заказ целиком \\
    \textbf{широкостолбцовые}   & запись на много узлов: события
      & соединения \\
    \textbf{графовые}           & «друзья друзей», кратчайший путь
      & таблицы и отчёты \\
    \textbf{поисковые}          & морфология, опечатки, ранжирование
      & источник правды \\
    \textbf{векторные}          & «похожее по смыслу»
      & точное совпадение \\
  \end{tabular}}

  \slidegap
  \aside{Каждое семейство --- ответ на вопрос, который реляционная база решает
    \lo{плохо}.}
\end{frame}

% 8
\begin{frame}{Четыре вопроса выбора}
  \slidetop
  \claim{Выбор базы --- решение про \hi{запросы}, а не про моду.}

  {\renewcommand{\arraystretch}{1.5}
  \begin{tabular}{@{}L{4.6cm}L{8.2cm}@{}}
    \textbf{какой формы данные}       & таблицы, документы, пары, сеть связей,
      измерения во времени \\
    \hi{\textbf{какие будут запросы}} & главный: что мы будем \hi{спрашивать} \\
    \textbf{какая согласованность}    & деньги или лайки \\
    \textbf{как масштабировать}       & одна большая машина или много узлов \\
  \end{tabular}}
\end{frame}

% 9
\begin{frame}{Одни данные, разная цена}
  \slidetop
  \claim{Заказы: номер, клиент, дата, сумма, позиции.}

  {\renewcommand{\arraystretch}{1.5}\small
  \begin{tabular}{@{}L{3.9cm}L{2.9cm}L{2.9cm}L{2.9cm}@{}}
    & \colhead{accent}{реляционная} & \colhead{alt}{документная}
      & \colhead{alt}{ключ-значение} \\
    заказ по номеру            & быстро & быстро & быстро \\
    сумма по клиентам за квартал & \hi{группировка по индексу} & обход всех
      документов & \lo{перебор всех ключей} \\
    заказы с товаром X         & \hi{соединение по индексу} & индекс по массиву,
      если есть & \lo{перебор всех ключей} \\
  \end{tabular}}

  \slidegap
  \aside{Стоимость вопросов различается в \lo{тысячи} раз --- и вы платите её
    каждый день.}
\end{frame}

% 10
\begin{frame}{Согласованность}
  \slidetop
  \versus{Транзакции}{%
    \begin{cbul}{accent}
      \item всё или ничего
      \item промежуточного никто не видит
      \item сохранил --- вижу, всегда
      \item для денег
    \end{cbul}}{Доступность}{%
    \begin{cbul}{alt}
      \item запись принял один узел
      \item остальные узнают потом
      \item сохранил --- \lo{не вижу}, секунды
      \item для лайков
    \end{cbul}}

  \slidegap
  \hrulethin{0.42}
  \rulegap
  \aside{Реплика для чтения отстаёт так же. Своё только что записанное читают
    оттуда, \hi{куда писали}.}
\end{frame}

% 11
\begin{frame}{Масштабирование}
  \slidetop
  \claim{Сначала \hi{измерить}, потом менять.}

  {\renewcommand{\arraystretch}{1.6}
  \begin{tabular}{@{}L{4.4cm}L{8.4cm}@{}}
    \hi{\textbf{реляционная}}    & одна большая машина и реплики для чтения ---
      на этом \hi{долго} \\
    \lo{\textbf{распределённая}} & много узлов по построению --- ценой
      согласованности \\
    \lo{\textbf{«на вырост»}}    & сложность каждый день за проблему, которой
      может не быть \\
  \end{tabular}}

  \slidegap
  \aside{Переезд, когда метрики показали предел, --- проект на месяцы, но с
    понятной причиной.}
\end{frame}

% 12
\begin{frame}[plain]
  \vote{Интернет-магазин: товары, заказы, клиенты, отчёты для бухгалтерии.
    Какая база?}
\end{frame}

% 13
\begin{frame}{Два ответа}
  \slidetop
  \versus{Магазин}{%
    \begin{cbul}{accent}
      \item отчёты поперёк заказов
      \item клиент в тысяче заказов
      \item запросы заранее не все известны
      \item \hi{реляционная}
    \end{cbul}}{Лента игры}{%
    \begin{cbul}{alt}
      \item миллион записей в секунду
      \item игрок --- ключ
      \item только «последние сто по игроку»
      \item \lo{широкостолбцовая}
    \end{cbul}}

  \slidegap
  \hrulethin{0.42}
  \rulegap
  \aside{Модный ответ выбирают по статье, правильный --- по запросам. Спросите:
    какие отчёты, по чему искать, что хранить как историю.}
\end{frame}

% 14
\begin{frame}{Базы сочетают}
  \slidetop
  \begin{center}
    \begin{tikzpicture}[font=\footnotesize]
      \node[box,draw=accent,line width=1.0pt,minimum width=3.6cm,
            minimum height=1.2cm] (pg) at (0,0)
        {\textbf{PostgreSQL}\\источник правды};
      \node[box,minimum width=3.6cm] (r) at (7.4,1.25)  {\textbf{Redis}\\кэш};
      \node[box,minimum width=3.6cm] (s) at (7.4,0)     {\textbf{поисковая}\\полнотекстовый поиск};
      \node[box,minimum width=3.6cm] (a) at (7.4,-1.25) {\textbf{аналитическая}\\отчёты};
      \draw[arr] (pg.east) -- (r.west);
      \draw[arr] (pg.east) -- node[above,font=\scriptsize,text=muted] {копия данных} (s.west);
      \draw[arr] (pg.east) -- (a.west);
    \end{tikzpicture}
  \end{center}

  \slidegap
  \begin{cbul}{alt}
    \item каждая база --- копии, обновления, мониторинг и человек, который её
      понимает
    \item транзакции между двумя базами \lo{не бывает}: согласованность ---
      на вас
  \end{cbul}
\end{frame}

% 15
\begin{frame}{Почему PostgreSQL}
  \slidetop
  \claim{Выбор \hi{по умолчанию}, если нет причины иначе.}

  \begin{cbul}{accent}
    \item открытый и бесплатный
    \item следует стандарту SQL строже большинства
    \item JSON, полнотекстовый поиск, расширения
    \item его знает любая инфраструктура и любое облако
  \end{cbul}

  \slidegap
  \aside{SQLite --- база в одном файле, для локальных утилит. В проде ---
    управляемый сервис: код тот же, копии и обновления делает провайдер.}
\end{frame}

% 16
\begin{frame}[plain]
  \keyline{Семейства отличаются тем, что они делают быстро.\\
    Базу выбирают по запросам, а не по моде.}
\end{frame}

% ==================================================== 2. RELATIONAL MODEL ==
\section{2. Реляционная модель}

% 17
\begin{frame}[plain]
  \sectionopener{2}{Реляционная модель}{Пять минут, чтобы \hi{договориться
    о словах}.}
\end{frame}

% 18
\begin{frame}[fragile]{Одна схема}
  \slidetop
  \begin{center}
    \begin{tikzpicture}[font=\footnotesize,
      tbl/.style={box,align=left,minimum width=3.9cm,inner sep=5pt}]
      \node[tbl] (c) at (0,0)
        {\textbf{customers}\\\ttfamily \hi{id} \ \ \ \ PK\\\ttfamily email UNIQUE\\\ttfamily name};
      \node[tbl] (o) at (7.6,0)
        {\textbf{orders}\\\ttfamily \hi{id} \ \ \ \ \ \ \ \ \ \ PK\\\ttfamily \lo{customer\_id} FK\\\ttfamily total};
      \draw[arrl] (o.west) -- node[above,font=\scriptsize,text=alt] {внешний ключ}
        node[below,font=\scriptsize,text=muted] {индекс по \texttt{customer\_id}} (c.east);
    \end{tikzpicture}
  \end{center}

\begin{lstlisting}[style=plain,basicstyle=\codesm]
select o.id, o.total, c.name
from orders o
join customers c on c.id = o.customer_id
where o.customer_id = 42;
\end{lstlisting}

  \rulegap
  \aside{Связь --- \lo{ссылкой}, а не вложенностью. SQL говорит «что»; «как» ---
    решает планировщик.}
\end{frame}

% 19
\begin{frame}{Словарь}
  \slidetop
  {\renewcommand{\arraystretch}{1.36}\small
  \begin{tabular}{@{}L{3.3cm}L{9.5cm}@{}}
    \hi{\textbf{ключи}}        & первичный определяет строку, внешний ссылается на
      неё \\
    \hi{\textbf{индекс}}       & находит строки, не читая таблицу; ускоряет
      чтение, \lo{замедляет} запись \\
    \hi{\textbf{план}}         & как база выполнит запрос: первый инструмент,
      когда медленно \\
    \hi{\textbf{соединение}}   & строки двух таблиц по условию --- тот же
      \texttt{Join} \\
    \hi{\textbf{транзакция}}   & все операции или ни одной; фиксация ждёт диск \\
    \hi{\textbf{нормализация}} & каждый факт хранится один раз \\
    \hi{\textbf{диалект}}      & свои типы и синтаксис поверх стандарта SQL \\
  \end{tabular}}

  \rulegap
  \aside{Запрос без индекса по большой таблице --- полное сканирование.}
\end{frame}

% 20
\begin{frame}{Ограничения}
  \slidetop
  \claim{Ограничения --- \hi{последняя линия}: их нельзя обойти.}

  \versus{Валидация в сервисе}{%
    \begin{cbul}{accent}
      \item проверяет запрос на границе
      \item только свой сервис
      \item два запроса сразу --- \lo{обойдут}
    \end{cbul}}{Ограничение в базе}{%
    \begin{cbul}{alt}
      \item проверяет данные при записи
      \item сервис, скрипт, другой сервис
      \item два запроса сразу --- один откажет
    \end{cbul}}

  \slidegap
  \aside{\texttt{not null}, \texttt{unique}, \texttt{check}, внешний ключ.
    Уникальность почты --- ограничение, а не только проверка в коде.}
\end{frame}

% ============================================================ 3. ADO.NET ==
\section{3. ADO.NET без магии}

% 21
\begin{frame}[plain]
  \sectionopener{3}{ADO.NET --- как оно работает без магии}{Под любой ORM лежит
    \hi{ровно это}.}
\end{frame}

% 22
\begin{frame}{Что под ORM}
  \slidetop
  \begin{center}
    \begin{tikzpicture}
      \foreach \i/\t in {1/{ваш код},2/{\textbf{ADO.NET}\\общие абстракции},
                         3/{\textbf{драйвер}\\Npgsql},4/{сетевой\\протокол},5/{база}} {
        \node[box,minimum width=2.55cm,minimum height=1.1cm,inner sep=2pt,
              font=\footnotesize] (l\i) at ({(\i-1)*2.85},0) {\t};
      }
      \foreach \i [evaluate=\i as \j using int(\i+1)] in {1,...,4}
        \draw[arr] (l\i) -- (l\j);
      \node[box,draw=accent,line width=1.0pt,minimum width=2.55cm,
            minimum height=1.1cm,inner sep=2pt,font=\footnotesize] at (2.85,0)
        {\textbf{ADO.NET}\\общие абстракции};
    \end{tikzpicture}
  \end{center}

  \slidegap
  \begin{cbul}{accent}
    \item пять абстракций: соединение, команда, параметр, читатель, транзакция
    \item драйвер реализует их и говорит с базой по её протоколу
    \item ORM внутри делает то, что сегодня сделаем \hi{руками}
  \end{cbul}
\end{frame}

% 23
\begin{frame}[plain]
  \vote{Открыть соединение, выполнить запрос, закрыть соединение. Что дорогого
    произошло при «открыть»?}
\end{frame}

% 24
\begin{frame}{Пул соединений}
  \slidetop
  \begin{center}
    \begin{tikzpicture}[font=\footnotesize]
      \node[box,minimum width=2.5cm,minimum height=1.9cm] (code) at (0,0) {ваш код};
      \node[draw=accent,line width=1.0pt,fill=codebg,minimum width=4.4cm,
            minimum height=1.9cm] (pool) at (5.9,0) {};
      \node[font=\scriptsize\bfseries,text=accent,above=0.25em] at (pool.north) {ПУЛ};
      \foreach \i in {1,...,4}
        \node[draw=hairline,line width=0.8pt,fill=paper,minimum width=0.7cm,
              minimum height=0.7cm] (s\i) at ({3.95+\i*0.78},0) {};
      \node[box,minimum width=2.3cm,minimum height=1.9cm] (db) at (11.5,0) {база};
      \draw[arra] ([yshift=0.45cm]code.east) -- node[above,font=\scriptsize,text=accent]
        {\texttt{Open}: взять} ([yshift=0.45cm]pool.west);
      \draw[arr] ([yshift=-0.45cm]pool.west) -- node[below,font=\scriptsize,text=muted]
        {\texttt{Close}: вернуть} ([yshift=-0.45cm]code.east);
      \draw[muted,line width=1.4pt] (pool.east) -- node[above,font=\scriptsize,text=muted]
        {TCP, TLS,} node[below,font=\scriptsize,text=muted] {вход} (db.west);
    \end{tikzpicture}
  \end{center}

  \slidegap
  {\large Настоящее соединение живёт \hi{долго}, ваш объект --- \lo{секунды}.}

  \rulegap
  \aside{Пул включён по умолчанию; предел --- порядка ста соединений на процесс,
    задаётся в строке подключения.}
\end{frame}

% 25
\begin{frame}{Пул конечен}
  \slidetop
  \claim{Пул --- на строку подключения и на \hi{процесс}.}

  {\renewcommand{\arraystretch}{1.6}
  \begin{tabular}{@{}L{4.6cm}L{8.2cm}@{}}
    \textbf{три экземпляра сервиса} & три пула: $3 \times 100 = 300$ соединений \\
    \lo{\textbf{предел базы}}       & свой, и у PostgreSQL по умолчанию невелик \\
    \hi{\textbf{сверяют при развёртывании}} & сумма пулов $\le$ предела базы \\
    \textbf{экземпляров много}      & внешний пул между сервисами и базой \\
  \end{tabular}}

  \slidegap
  \aside{При возврате драйвер сбрасывает транзакцию и параметры сессии: каждое
    открытие --- \hi{чистое} соединение.}
\end{frame}

% 26
\begin{frame}{Два правила}
  \slidetop
  \versus{Открывать поздно, закрывать рано}{%
    \begin{cbul}{accent}
      \item взять ровно перед запросом
      \item вернуть сразу после
      \item не держать весь HTTP-запрос
    \end{cbul}}{Всегда using}{%
    \begin{cbul}{alt}
      \item \texttt{Dispose} --- это возврат в пул
      \item забыли --- не вернётся
      \item сто забытых --- пул пуст
    \end{cbul}}

  \slidegap
  \hrulethin{0.42}
  \rulegap
  \aside{Такси вызывают на поездку и отпускают по приезде. Машину у подъезда
    «на всякий случай» отнимают у ста других вызовов.}
\end{frame}

% 27
\begin{frame}[fragile]{Утечка}
  \slidetop
\begin{lstlisting}[style=csharp,basicstyle=\codesm]
for (int i = 1; ; i++)
{
    var c = await small.OpenConnectionAsync();   // no using
    Console.WriteLine($"opened {i}");
}
\end{lstlisting}

\begin{lstlisting}[style=term,basicstyle=\termsm]
$ dotnet run -- leak
opened 1
opened 2
opened 3
NpgsqlException: The connection pool has been exhausted, either raise 'Max Pool Size' (currently 3)
or 'Timeout' (currently 3 seconds) in your connection string.
\end{lstlisting}

  \rulegap
  \aside{Выглядит как «база тормозит», хотя база \lo{свободна}. Метрика пула
    «занято» на уровне предела --- подпись утечки.}
\end{frame}

% 28
\begin{frame}[fragile]{Одно значение}
  \vspace{0.3em}
\begin{lstlisting}[style=csharp,basicstyle=\codexs,escapechar=|]
|\mk{c1}|public sealed class OrderCounter(NpgsqlDataSource dataSource)
{
    public async Task<long> CountForCustomerAsync(
        int customerId, CancellationToken ct)
    {
|\mk{c2}|        await using var connection =
            await dataSource.OpenConnectionAsync(ct);
        await using var command = new NpgsqlCommand(
            "select count(*) from orders " +
|\mk{c3}|            "where customer_id = @customerId", connection);
        command.Parameters.AddWithValue("customerId", customerId);

|\mk{c4}|        var result = await command.ExecuteScalarAsync(ct);
        return (long)result!;
    }
}
\end{lstlisting}
  \begin{tikzpicture}[remember picture,overlay,every node/.style={rawlab}]
    \node at ([xshift=10.3cm]c1) {источник --- singleton};
    \node at ([xshift=10.3cm]c2) {открыть перед запросом};
    \node[text=accent] at ([xshift=10.3cm]c3) {параметр, не склейка};
    \node at ([xshift=10.3cm]c4) {одно значение};
  \end{tikzpicture}

  \aside{\texttt{await using} возвращает соединение в пул --- при \hi{любом}
    выходе.}
\end{frame}

% 29
\begin{frame}[fragile]{Что стоит за строками}
  \slidetop
  {\renewcommand{\arraystretch}{1.4}\small
  \begin{tabular}{@{}L{3.9cm}L{8.9cm}@{}}
    \hi{\textbf{источник данных}}   & один на приложение, владеет пулом \\
    \hi{\textbf{строка подключения}} & адрес, база, пользователь, пул; с паролем
      --- \lo{секрет} \\
    \hi{\textbf{токен}}             & клиент ушёл --- драйвер шлёт базе отмену
      запроса \\
    \hi{\textbf{тип параметра}}     & выводится из значения; \texttt{null} ---
      это \texttt{DBNull.Value} \\
    \hi{\textbf{\texttt{(long)}}}   & \texttt{count} --- восьмибайтовое целое \\
  \end{tabular}}

  \rulegap
\begin{lstlisting}[style=term,basicstyle=\termsm]
$ dotnet run -- count
2
InvalidCastException: Unable to cast object of type 'System.Int64' to type 'System.Int32'.
\end{lstlisting}

  \aside{С \texttt{(int)} --- исключение при выполнении, а \lo{не} ошибка
    компиляции.}
\end{frame}

% 30
\begin{frame}{Три вида выполнения}
  \slidetop
  \claim{Выбор метода --- это выбор \hi{формы результата}.}

  {\renewcommand{\arraystretch}{1.6}
  \begin{tabular}{@{}L{4.4cm}L{3.4cm}L{5.0cm}@{}}
    \texttt{ExecuteScalarAsync}   & одно значение   & \texttt{count}, \texttt{max},
      \texttt{returning id} \\
    \hi{\texttt{ExecuteReaderAsync}} & \hi{поток строк} & всё, что возвращает
      строки \\
    \texttt{ExecuteNonQueryAsync} & сколько строк затронуто & \texttt{insert},
      \texttt{update}, \texttt{delete} \\
  \end{tabular}}

  \slidegap
  \aside{\texttt{update} по идентификатору затронул \lo{ноль} строк --- это «не
    найдено». Синхронных версий в сервисе нет ни одной.}
\end{frame}

% 31
\begin{frame}[fragile]{Два null}
  \slidetop
  \claim{Строк нет --- \texttt{null}. Строка есть, значение пустое ---
    \lo{\texttt{DBNull}}.}

\begin{lstlisting}[style=plain,basicstyle=\codesm]
select max(shipped_at) from orders where customer_id = 2   -- no shipped orders
select shipped_at from orders where customer_id = 99       -- no such customer
\end{lstlisting}

\begin{lstlisting}[style=term,basicstyle=\termsm]
$ dotnet run -- nulls
max, no shipped: DBNull
no rows: null
\end{lstlisting}

  \rulegap
  \aside{Приведение к числу упадёт на обоих, по-разному. \texttt{count} всегда
    возвращает строку; «максимальную дату» проверяют \hi{явно}.}
\end{frame}

% 32
\begin{frame}{Число поездок}
  \slidetop
  \claim{Число поездок к базе важнее размера каждой.}

  \versus{Один запрос}{%
    \begin{cbul}{accent}
      \item заказы --- одним запросом
      \item позиции --- соединением
      \item \hi{2} поездки
    \end{cbul}}{Запрос в цикле}{%
    \begin{cbul}{alt}
      \item сто заказов
      \item для каждого --- запрос за позициями
      \item \lo{101} поездка
    \end{cbul}}

  \slidegap
  \hrulethin{0.42}
  \rulegap
  \aside{Проблема N + 1 живёт в любом коде, который запрашивает в цикле, а не
    только в ORM.}
\end{frame}

% 33
\begin{frame}[fragile]{Страница}
  \slidetop
  \versus{По ключу}{%
    \begin{cbul}{accent}
      \item «следующие после этого ключа»
      \item по индексу
      \item любая страница --- за одно время
    \end{cbul}}{Пропустить и взять}{%
    \begin{cbul}{alt}
      \item «пропусти 20\,000, дай 20»
      \item база перебирает пропущенные
      \item тысячная страница \lo{медленнее}
    \end{cbul}}

  \slidegap
\begin{lstlisting}[style=plain,basicstyle=\codesm]
select id, total from orders where id > @lastId order by id limit 20;
select id, total from orders order by id offset 20000 limit 20;
\end{lstlisting}

  \rulegap
  \aside{В контракте API это курсор вместо номера страницы.}
\end{frame}

% 34
\begin{frame}[fragile]{Читатель}
  \vspace{0.3em}
\begin{lstlisting}[style=csharp,basicstyle=\codexs,escapechar=|]
|\mk{r1}|await using var reader = await command.ExecuteReaderAsync(ct);
|\mk{r2}|int idOrdinal = reader.GetOrdinal("id"),
    totalOrdinal = reader.GetOrdinal("total"),
    shippedOrdinal = reader.GetOrdinal("shipped_at");

|\mk{r3}|while (await reader.ReadAsync(ct))
{
|\mk{r4}|    yield return new OrderRow(
        reader.GetInt32(idOrdinal),
        reader.GetDecimal(totalOrdinal),
|\mk{r5}|        await reader.IsDBNullAsync(shippedOrdinal, ct)
            ? null : reader.GetDateTime(shippedOrdinal));
}

public sealed record OrderRow(int Id, decimal Total, DateTime? ShippedAt);
\end{lstlisting}
  \begin{tikzpicture}[remember picture,overlay,every node/.style={rawlab}]
    \node[text=accent] at ([xshift=9.6cm]r1) {читатель --- поток};
    \node at ([xshift=9.6cm]r2) {номера --- один раз, до цикла};
    \node at ([xshift=9.6cm]r3) {строки приходят по одной};
    \node at ([xshift=9.6cm]r4) {наружу --- по мере чтения};
    \node at ([xshift=9.6cm]r5) {проверка \texttt{null}};
  \end{tikzpicture}

  \aside{Тип не совпал --- исключение при чтении: компилятор схему
    \lo{не видит}.}
\end{frame}

% 35
\begin{frame}[fragile]{Поток строк}
  \slidetop
  \claim{Не список и не массив: в памяти --- \hi{одна} строка.}

\begin{lstlisting}[style=term,basicstyle=\termsm]
$ dotnet run -- read
OrderRow { Id = 1, Total = 300.00, ShippedAt = 11/10/2026 9:30:00 AM }
OrderRow { Id = 2, Total = 120.00, ShippedAt =  }
\end{lstlisting}

  \slidegap
  \begin{cbul}{accent}
    \item тот же контракт, что перечислитель: курсор перед первой строкой
    \item каждое \texttt{ReadAsync} --- ожидание сети, поток не занят
    \item от базы до HTTP-ответа --- без единого списка в памяти
    \item соединение открыто, пока последовательность \hi{перебирают}
  \end{cbul}
\end{frame}

% 36
\begin{frame}[plain]
  \vote{Вы прочитали три строки из читателя и вышли из цикла по
    \texttt{break}. Что с соединением?}
\end{frame}

% 37
\begin{frame}[fragile]{Один читатель на соединение}
  \slidetop
  \claim{Пока читатель открыт, соединение \lo{занято} им.}

\begin{lstlisting}[style=csharp,basicstyle=\codesm]
var reader = await cmd.ExecuteReaderAsync();          // no using
await reader.ReadAsync();
await using var cmd2 = new NpgsqlCommand("select count(*) from orders", c);
await cmd2.ExecuteScalarAsync();
\end{lstlisting}

\begin{lstlisting}[style=term,basicstyle=\termsm]
$ dotnet run -- break
NpgsqlOperationInProgressException: A command is already in progress: select id from orders
\end{lstlisting}

  \rulegap
  \aside{\texttt{using} закрывает читатель и дочитывает остаток. Он не
    украшение: он делает \texttt{break} \hi{безопасным}.}
\end{frame}

% 38
\begin{frame}{Общие типы}
  \slidetop
  \claim{Переносимость \hi{механики}, не переносимость SQL.}

  \versus{Переносится}{%
    \begin{cbul}{accent}
      \item \texttt{DbConnection}, \texttt{DbCommand}
      \item код не знает, какая под ним база
      \item переключается регистрацией
    \end{cbul}}{Не переносится}{%
    \begin{cbul}{alt}
      \item \texttt{returning}, типы дат, регистр имён
      \item диалект остаётся в тексте запроса
      \item тесты на SQLite, прод на PostgreSQL
    \end{cbul}}

  \slidegap
  \hrulethin{0.42}
  \rulegap
  \aside{Тесты доступа к данным гоняют на той базе, что в проде, --- в
    контейнере.}
\end{frame}

% 39
\begin{frame}[fragile]{Транзакция: вставка}
  \slidetop
  \claim{Создать заказ и уменьшить остаток: вторая упала --- первой \lo{нет}.}

\begin{lstlisting}[style=csharp,basicstyle=\codexs,escapechar=|]
await using var connection = await dataSource.OpenConnectionAsync(ct);
|\mk{x1}|await using var transaction =
    await connection.BeginTransactionAsync(ct);

await using var insert = new NpgsqlCommand(
    "insert into orders (customer_id, product_id, quantity) " +
|\mk{x2}|    "values (@c, @p, @q) returning id", connection, transaction);
insert.Parameters.AddWithValue("c", customerId);
insert.Parameters.AddWithValue("p", productId);
insert.Parameters.AddWithValue("q", quantity);
|\mk{x3}|var orderId = (int)(await insert.ExecuteScalarAsync(ct))!;
\end{lstlisting}
  \begin{tikzpicture}[remember picture,overlay,every node/.style={rawlab}]
    \node[text=accent] at ([xshift=9.9cm]x1) {начать на соединении};
    \node at ([xshift=9.9cm]x2) {команда в транзакции};
    \node at ([xshift=9.9cm]x3) {идентификатор --- сразу};
  \end{tikzpicture}

  \slidegap
  \aside{\texttt{returning} --- диалект PostgreSQL: вставка и идентификатор одной
    поездкой.}
\end{frame}

% 40
\begin{frame}[fragile]{Транзакция: условие и фиксация}
  \vspace{0.3em}
\begin{lstlisting}[style=csharp,basicstyle=\codexs,escapechar=|]
await using var update = new NpgsqlCommand(
    "update products set stock = stock - @q " +
|\mk{y1}|    "where id = @p and stock >= @q", connection, transaction);
update.Parameters.AddWithValue("q", quantity);
update.Parameters.AddWithValue("p", productId);
|\mk{y2}|if (await update.ExecuteNonQueryAsync(ct) == 0)
    throw new InsufficientStockException(productId);

|\mk{y3}|await transaction.CommitAsync(ct);
return orderId;
\end{lstlisting}
  \begin{tikzpicture}[remember picture,overlay,every node/.style={rawlab}]
    \node at ([xshift=9.9cm]y1) {проверка --- в самом SQL};
    \node at ([xshift=9.9cm]y2) {ноль строк --- остатка нет};
    \node[text=accent] at ([xshift=9.9cm]y3) {зафиксировать в конце};
  \end{tikzpicture}
  \vspace{-1.2em}

\begin{lstlisting}[style=term,basicstyle=\termsm]
$ dotnet run -- tx
order 4
InsufficientStockException: Insufficient stock for product 1
orders: 4, stock: 0
\end{lstlisting}

  \aside{Явного отката \lo{нет}: вышли без фиксации --- драйвер откатил сам.}
\end{frame}

% 41
\begin{frame}{Проверка --- в одной команде}
  \slidetop
  \versus{Условие в update}{%
    \begin{cbul}{accent}
      \item проверка и изменение неразделимы
      \item второй запрос видит уже ноль
      \item затронуто ноль строк --- отказ
    \end{cbul}}{Прочитать, проверить, записать}{%
    \begin{cbul}{alt}
      \item три шага из кода
      \item оба видят остаток «один»
      \item продано \lo{два} при одном на складе
    \end{cbul}}

  \slidegap
  \hrulethin{0.42}
  \rulegap
  {\large Что должно быть атомарным --- делается \hi{одной командой} базы.}
\end{frame}

% 42
\begin{frame}{Правила транзакции}
  \slidetop
  {\renewcommand{\arraystretch}{1.45}\small
  \begin{tabular}{@{}L{3.9cm}L{8.9cm}@{}}
    \hi{\textbf{короткая}}        & внутри только команды к базе; вызовы наружу
      --- \hi{до} неё \\
    \hi{\textbf{откат сам}}       & через \texttt{using}; забыли \texttt{CommitAsync} ---
      откатится \\
    \hi{\textbf{вложенных нет}}   & есть точки сохранения, и нужны они редко \\
    \hi{\textbf{повтор}}          & временная ошибка --- по коду состояния; вся
      транзакция с начала \\
    \hi{\textbf{версия строки}}   & «обновить, где версия та же»: ноль строк ---
      \lo{409}, перечитать \\
    \hi{\textbf{блокировка при чтении}} & второй ждёт первого: для частых
      конфликтов \\
  \end{tabular}}

  \rulegap
  \aside{Пока транзакция открыта, база держит блокировки на изменённых строках
    --- остальные \lo{стоят}.}
\end{frame}

% 43
\begin{frame}[plain]
  \vote{Запрос собран плюсом: текст SQL и значение из строки запроса. Что не
    так?}
\end{frame}

% 44
\begin{frame}[fragile]{Инъекция}
  \slidetop
  \claim{Данные попали в \lo{текст программы}.}

\begin{lstlisting}[style=csharp,basicstyle=\codesm]
var sql = "select count(*) from orders where customer_id = " + input;
\end{lstlisting}

\begin{lstlisting}[style=term,basicstyle=\termsm]
$ dotnet run -- inject
concat [1] -> 2
concat [1 or 1=1] -> 3
concat [1; drop table orders; --] -> 2
PostgresException: 42P01: relation "orders" does not exist
\end{lstlisting}

  \rulegap
  \begin{cbul}{alt}
    \item \texttt{or 1=1} --- все заказы всех клиентов: тихая выгрузка
    \item \texttt{drop table} --- второй запрос: таблицы \lo{нет}
  \end{cbul}
\end{frame}

% 45
\begin{frame}[fragile]{Параметр}
  \slidetop
  \claim{Значение уходит \hi{отдельно} от текста.}

\begin{lstlisting}[style=csharp,basicstyle=\codesm]
var sql = "select count(*) from orders where customer_id::text = @customerId";
cmd.Parameters.AddWithValue("customerId", input);
\end{lstlisting}

\begin{lstlisting}[style=term,basicstyle=\termsm]
$ dotnet run -- param
param [1] -> 2
param [1 or 1=1] -> 0
param [1; drop table orders; --] -> 0
3
\end{lstlisting}

  \begin{cbul}{accent}
    \item база не разбирает значение как SQL: оно \hi{не в SQL}
    \item экранирование руками --- гонка с теми, кто ищет обходы
    \item один текст --- один план на все значения: ещё и быстрее
  \end{cbul}
\end{frame}

% 46
\begin{frame}[plain]
  \statement{В текст запроса не попадает \lo{ни один байт}, пришедший снаружи.
    Только \hi{параметры}.}
\end{frame}

% 47
\begin{frame}{Где параметр не спасает}
  \slidetop
  {\renewcommand{\arraystretch}{1.6}
  \begin{tabular}{@{}L{4.4cm}L{8.4cm}@{}}
    \lo{\textbf{шаблон \texttt{like}}}    & \texttt{\%} в значении найдёт всё:
      экранировать или запретить \\
    \lo{\textbf{имена таблиц и колонок}}  & не параметризуются: \hi{белый список}
      констант в коде \\
    \lo{\textbf{сырой SQL в ORM}}         & склеили строку --- инъекция
      вернулась \\
  \end{tabular}}

  \slidegap
  \aside{«Снаружи» --- это и строка запроса, и тело, и заголовки, и сама база: в
    неё всё тоже когда-то пришло снаружи.}
\end{frame}

% 48
\begin{frame}[plain]
  \vote{Что не так со звёздочкой в \texttt{select *}?}
\end{frame}

% 49
\begin{frame}{Колонки --- явно}
  \slidetop
  \claim{Колонки перечисляются \hi{всегда}.}

  \begin{cbul}{alt}
    \item читаете ненужное: трафик и память на каждой строке
    \item новая колонка \lo{сдвигает} порядковые номера
    \item в коде не видно, от каких колонок он зависит
  \end{cbul}

  \slidegap
  \hrulethin{0.42}
  \rulegap
  \aside{SQL живёт в классе хранилища, рядом с методом: его читают и ревьюят как
    код. Ошибки базы там же становятся исключениями домена.}
\end{frame}

% 50
\begin{frame}[plain]
  \keyline{Соединение: взять поздно, вернуть рано.\\
    Читатель --- поток. Транзакция откатывается сама.\\
    Данные не должны становиться программой.}
\end{frame}

% ======================================================= 4. BRIDGE TO ORM ==
\section{4. Мост к ORM}

% 51
\begin{frame}[plain]
  \sectionopener{4}{Мост к ORM}{ORM не заменяет драйвер: она пишет то, что вы
    писали \hi{руками}.}
\end{frame}

% 52
\begin{frame}[plain]
  \vote{В листинге с читателем: сколько строк про то, что мы хотим получить, и
    сколько про то, как?}
\end{frame}

% 53
\begin{frame}{Три боли}
  \slidetop
  {\renewcommand{\arraystretch}{1.7}
  \begin{tabular}{@{}L{3.8cm}L{9.0cm}@{}}
    \lo{\textbf{маппинг руками}}  & десять строк механики на строку смысла;
      ошибка --- при чтении \\
    \lo{\textbf{дублирование}}    & тот же запрос в пяти местах; строку SQL
      дополняют только склейкой \\
    \lo{\textbf{схема и типы}}    & колонку переименовали --- компилятор
      \lo{не заметил} \\
  \end{tabular}}

  \slidegap
  \hrulethin{0.42}
  \rulegap
  \aside{В списке нет «SQL сложный». 20 таблиц $\times$ 10 запросов $\times$ 10
    строк --- две тысячи строк без единого бизнес-решения. Для пяти запросов ORM
    \hi{не нужна}.}
\end{frame}

% 54
\begin{frame}{Алгебра и LINQ}
  \slidetop
  \claim{Те же операции --- в другом синтаксисе. Это \hi{дизайн}.}

  {\renewcommand{\arraystretch}{1.25}\small
  \begin{tabular}{@{}L{4.2cm}L{4.2cm}L{4.2cm}@{}}
    \eyebrow{реляционная алгебра} & \colhead{alt}{SQL} & \colhead{accent}{LINQ} \\
    выборка     & \texttt{where}    & \texttt{Where} \\
    проекция    & \texttt{select}   & \texttt{Select} \\
    соединение  & \texttt{join}     & \texttt{Join} \\
    группировка & \texttt{group by} & \texttt{GroupBy} \\
    объединение & \texttt{union}    & \texttt{Union} \\
    сортировка  & \texttt{order by} & \texttt{OrderBy} \\
  \end{tabular}}

  \rulegap
  \aside{Края есть: оконные функции, рекурсивные запросы. Запрос на LINQ стал
    длиннее своего SQL --- пишите SQL.}
\end{frame}

% 55
\begin{frame}[fragile]{Из LINQ в SQL}
  \vspace{0.3em}
\begin{lstlisting}[style=csharp,basicstyle=\codesm,escapechar=`]
var paid = context.Orders
`\mk{q1}`    .Where(o => o.CustomerId == customerId
             && o.Status == OrderStatus.Paid)
`\mk{q2}`    .OrderByDescending(o => o.CreatedAt)
`\mk{q3}`    .Select(o => new { o.Id, o.Total })
`\mk{q4}`    .Take(10);
\end{lstlisting}
  \begin{tikzpicture}[remember picture,overlay,every node/.style={rawlab}]
    \node at ([xshift=9.2cm]q1) {\texttt{WHERE}};
    \node at ([xshift=9.2cm]q2) {\texttt{ORDER BY ... DESC}};
    \node[text=accent] at ([xshift=9.2cm]q3) {\texttt{SELECT}: две колонки};
    \node at ([xshift=9.2cm]q4) {\texttt{LIMIT}};
  \end{tikzpicture}
  \vspace{-1.2em}

\begin{lstlisting}[style=plain,basicstyle=\codesm]
SELECT o.id, o.total
FROM orders AS o
WHERE o.customer_id = @customerId AND o.status = 'Paid'
ORDER BY o.created_at DESC
LIMIT 10
\end{lstlisting}

  \rulegap
  \aside{\texttt{customerId} стал \hi{параметром}: провайдер не вклеивает
    значения. Инъекция невозможна по построению.}
\end{frame}

% 56
\begin{frame}{Механизм}
  \slidetop
  \begin{center}
    \begin{tikzpicture}
      \foreach \i/\t in {1/лямбда,2/{дерево\\выражений},3/{провайдер\\обходит},
                         4/{SQL и\\параметры},5/{драйвер\\ADO.NET},6/читатель,
                         7/объекты} {
        \node[box,minimum width=1.62cm,minimum height=1.1cm,inner sep=1pt,
              font=\scriptsize] (m\i) at ({(\i-1)*1.95},0) {\t};
      }
      \foreach \i [evaluate=\i as \j using int(\i+1)] in {1,...,6}
        \draw[arr] (m\i) -- (m\j);
      \node[box,draw=accent,line width=1.0pt,minimum width=1.62cm,
            minimum height=1.1cm,inner sep=1pt,font=\scriptsize\bfseries]
        at (7.8,0) {драйвер\\ADO.NET};
      \draw[rawbrace] ([yshift=-0.25cm]m4.south west) -- ([yshift=-0.25cm]m1.south west);
      \draw[rawbrace] ([yshift=-0.25cm]m7.south east) -- ([yshift=-0.25cm]m5.south west);
      \node[font=\scriptsize,text=muted] at ([yshift=-0.75cm]$(m1.south west)!0.5!(m4.south west)$)
        {пишет ORM};
      \node[font=\scriptsize,text=accent] at ([yshift=-0.75cm]$(m5.south west)!0.5!(m7.south east)$)
        {сегодня --- руками};
    \end{tikzpicture}
  \end{center}

  \slidegap
  \begin{cbul}{accent}
    \item ORM --- генератор SQL и генератор маппинга \hi{поверх драйвера}
    \item перевод кэшируется: микросекунды против миллисекунд сети
    \item цена: абстракция \lo{протекает}, когда вы перестали смотреть на SQL
  \end{cbul}
\end{frame}

% 57
\begin{frame}[plain]
  \vote{В условие \texttt{Where} вы написали вызов своего метода. Над списком
    это работало. Что произойдёт над таблицей?}
\end{frame}

% 58
\begin{frame}{Оба варианта хуже}
  \slidetop
  \versus{Исключение}{%
    \begin{cbul}{accent}
      \item «не удалось перевести»
      \item честно, но при выполнении
      \item тесты над списком не поймают
    \end{cbul}}{Вычисление на клиенте}{%
    \begin{cbul}{alt}
      \item в базу уходит запрос без условия
      \item \lo{миллион} строк в память ради десяти
      \item в логе SQL выглядит прилично
    \end{cbul}}

  \slidegap
  \hrulethin{0.42}
  \rulegap
  \aside{Родственный сюрприз: сто заказов, покупатель каждого по обращению ---
    сто запросов. Оба лечатся одним: \hi{смотреть SQL}.}
\end{frame}

% 59
\begin{frame}{IQueryable --- это описание SQL}
  \slidetop
  \claim{Три шага, минута на запрос.}

  {\renewcommand{\arraystretch}{1.6}
  \begin{tabular}{@{}L{1.0cm}L{11.8cm}@{}}
    \hi{\textbf{1}} & представьте SQL: какие таблицы, условия, колонки \\
    \hi{\textbf{2}} & напишите LINQ так, чтобы он лёг на этот SQL \\
    \hi{\textbf{3}} & посмотрите в лог: \hi{совпало} ли \\
  \end{tabular}}

  \slidegap
  \aside{ORM --- инструмент, а не архитектура: DTO, интерфейсы хранилищ и
    границы остаются в силе.}
\end{frame}

% 60
\begin{frame}[plain]
  \keyline{ORM --- генератор SQL и маппинга поверх драйвера.\\
    Граница между базой и памятью --- ваша забота.}
\end{frame}

\end{document}
